\subsection{Diagonalization and Simultaneous Forms}

\begin{perspective} \label{per:diagonalizability-viewpoint} \hlt{Diagonalization as Exhaustion by Eigenspaces}\\
Distinct eigenspaces are already separated by Proposition \ref{prop:eigenspaces-for-distinct-eigenvalues-form-direct-sum}. Thus diagonalizability is the stronger statement that these separated pieces account for every direction in $\mathcal{V}$.
\end{perspective}

\begin{theorem} \label{thm:algebraic-spectral-resolution-diagonalizable} \hlt{Algebraic Spectral Resolution}\\
Let $0\neq\mathcal{V}$ be finite-dimensional, let $T\in\LL(\mathcal{V})$, and let $\Lambda=\{\lambda_1,\ldots,\lambda_m\}\subseteq\mathbb{F}$ with $\lambda_i\neq\lambda_j$ for $i\neq j$. The following are equivalent:
\begin{enumerate}[label=(\roman*).]
\setlength{\itemsep}{0pt}
\item $T$ is diagonalizable and $\operatorname{spec}(T)=\Lambda$.
\item $\exists$ nonzero projections $\rho_1,\ldots,\rho_m\in\LL(\mathcal{V})$ s.t. $\rho_i\rho_j=0\ (i\neq j)$, $\sum\limits_{i=1}^{m}\rho_i=I_{\mathcal{V}}$, $T=\sum\limits_{i=1}^{m}\lambda_i\rho_i$.
\end{enumerate}
In this case, for each $i$, $\operatorname{range}(\rho_i)=E_{\lambda_i}(T)$; $\ker(\rho_i)=\bigoplus_{j\neq i}E_{\lambda_j}(T)$; and $\rho_i=\prod\limits_{j\neq i}\frac{T-\lambda_jI_{\mathcal{V}}}{\lambda_i-\lambda_j}$
\begin{proof}
\begin{enumerate}[label=(\roman*).]
\setlength{\itemsep}{0pt}
\item[$(i)$]$\Rightarrow (ii)$ By Theorem \ref{thm:conditions-equivalent-to-diagonalizability}, $\mathcal{V}=\bigoplus\limits_{i=1}^{m}E_{\lambda_i}(T)$. Let $\rho_i=\rho_{E_{\lambda_i}(T),\,\bigoplus_{j\neq i}E_{\lambda_j}(T)}$.\\
Then $\rho_i\rho_j=0$ for $i\neq j$, $\sum\limits_{i=1}^{m}\rho_i=I_{\mathcal{V}}$, and $T=\sum\limits_{i=1}^{m}\lambda_i\rho_i$. Also $\operatorname{range}(\rho_i)$ and $\ker(\rho_i)$ are as displayed.\\
For $\ell_i(z)=\prod\limits_{j\neq i}\frac{z-\lambda_j}{\lambda_i-\lambda_j}$, we have $\ell_i(\lambda_k)=\delta_{ik}$, hence $\ell_i(T)=\rho_i$.
\item[$(ii)$]$\Rightarrow (i)$ Put $\mathcal{V}_i=\operatorname{range}(\rho_i)$. By Theorem \ref{thm:resolution-of-identity-gives-direct-sum}, $\mathcal{V}=\bigoplus\limits_{i=1}^{m}\mathcal{V}_i$.\\
If $\mathbf{x}\in\mathcal{V}_i$, then $\rho_i\mathbf{x}=\mathbf{x}$ and $\rho_j\mathbf{x}=\mathbf{0}$ for $j\neq i$, so $T\mathbf{x}=\lambda_i\mathbf{x}$. Thus $\mathcal{V}_i\subseteq E_{\lambda_i}(T)$.\\
Conversely, if $T\mathbf{x}=\lambda_i\mathbf{x}$ and $\mathbf{x}=\sum\limits_{j=1}^{m}\mathbf{x}_j$ with $\mathbf{x}_j\in\mathcal{V}_j$, then $\sum\limits_{j=1}^{m}(\lambda_j-\lambda_i)\mathbf{x}_j=\mathbf{0}$, so $\mathbf{x}_j=\mathbf{0}$ for $j\neq i$.\\
Hence $\mathcal{V}_i=E_{\lambda_i}(T)$. Since each $\mathcal{V}_i\neq0$, each $\lambda_i\in\operatorname{spec}(T)$; concatenating bases of the $\mathcal{V}_i$ gives an eigenbasis, so $T$ is diagonalizable and $\operatorname{spec}(T)=\Lambda$.\\
The kernel and polynomial formulas follow from $\mathcal{V}=\bigoplus\limits_{i=1}^{m}\mathcal{V}_i$ and $\ell_i(\lambda_k)=\delta_{ik}$.
\end{enumerate}
\hlt{Structural Commentary}: A diagonalizable operator is a weighted sum of the projections onto its eigenspaces.\\
\hlt{Key Insight}: Spectral projections are Lagrange polynomials in the operator.
\end{proof}
\end{theorem}

\begin{lemma} \label{lem:diagonalizable-restriction-to-invariant-subspace} \hlt{Diagonalizable Restrictions to Invariant Subspaces}\\
Let $T\in\LL(\mathcal{V})$ be diagonalizable and let $\mathcal{U}\in\operatorname{Inv}(T)$. For $\lambda\in\operatorname{spec}(T)$, set $\mathcal{U}_{\lambda}=\mathcal{U}\cap E_{\lambda}(T)$.\\
Then $\mathcal{U}=\bigoplus\limits_{\lambda\in\operatorname{spec}(T)}\mathcal{U}_{\lambda}$, where $T|_{\mathcal{U}}$ is diagonalizable.
\begin{proof}
Let $\rho_{\lambda}=p_{\lambda}(T)$ be the spectral projection from Theorem \ref{thm:algebraic-spectral-resolution-diagonalizable}. Since $\mathcal{U}\in\operatorname{Inv}(T)$, $\rho_{\lambda}(\mathcal{U})\subseteq\mathcal{U}$.\\
Hence, for $\mathbf{u}\in\mathcal{U}$, $\mathbf{u}=\sum\limits_{\lambda\in\operatorname{spec}(T)}\rho_{\lambda}\mathbf{u}$, where $\rho_{\lambda}\mathbf{u}\in\mathcal{U}_{\lambda}$. Thus $\mathcal{U}=\sum\limits_{\lambda\in\operatorname{spec}(T)}\mathcal{U}_{\lambda}$, and the sum is direct since $\mathcal{U}_{\lambda}\subseteq E_{\lambda}(T)$. Bases of the nonzero $\mathcal{U}_{\lambda}$ concatenate to an eigenbasis for $T|_{\mathcal{U}}$.\\
\hlt{Structural Commentary}: Invariant subspaces inherit spectral projections.\\
\hlt{Key Insight}: Polynomial spectral projections cannot move vectors out of a $T$-invariant subspace.
\end{proof}
\end{lemma}

\begin{theorem} \label{thm:simultaneous-diagonalization-commuting-diagonalizable-family} \hlt{Simultaneous Diagonalisation of Commuting Diagonalizable Operators}\\
Let $\dim\mathcal{V}<\infty$, and let $\emptyset\neq\mathcal{F}\subseteq\LL(\mathcal{V})$ be finite with every $T\in\mathcal{F}$ diagonalizable.\\
Then $\forall S,T\in\mathcal{F},\ ST=TS$ $\Longleftrightarrow$ $\exists$ a basis $\beta$ of $\mathcal{V}$ such that $[T]_{\beta\leftarrow\beta}$ is diagonal $\forall T\in\mathcal{F}$.
\begin{proof}
\begin{enumerate}[label=\roman*.]
\setlength{\itemsep}{0pt}
\item[$(\Rightarrow)$] Induct on $n=\dim\mathcal{V}$. If every $T\in\mathcal{F}$ is scalar, any basis works. Otherwise choose nonscalar $S\in\mathcal{F}$.\\
By Theorem \ref{thm:algebraic-spectral-resolution-diagonalizable}, $\mathcal{V}=\bigoplus\limits_{\lambda\in\operatorname{spec}(S)}E_{\lambda}(S)$, with at least two nonzero summands of dimension $<n$.\\
If $T\in\mathcal{F}$ and $\mathbf{x}\in E_{\lambda}(S)$, then $S(T\mathbf{x})=T(S\mathbf{x})=\lambda T\mathbf{x}$, so $E_{\lambda}(S)\in\operatorname{Inv}(T)$.\\
By Lemma \ref{lem:diagonalizable-restriction-to-invariant-subspace}, each $T|_{E_{\lambda}(S)}$ is diagonalizable; the restricted family still commutes.\\
Induction gives a common eigenbasis of each $E_{\lambda}(S)$; concatenate over $\lambda$.
\item[$(\Leftarrow)$] For $S,T\in\mathcal{F}$, set $D_S=[S]_{\beta\leftarrow\beta}$ and $D_T=[T]_{\beta\leftarrow\beta}$.\\
Since $D_S,D_T$ are diagonal, $[ST]_{\beta\leftarrow\beta}=D_SD_T=D_TD_S=[TS]_{\beta\leftarrow\beta}$. Thus $ST=TS$.
\end{enumerate}
\hlt{Structural Commentary}: Commutation preserves eigenspaces.\\
\hlt{Key Insight}: Diagonalize one operator, then restrict.
\end{proof}
\end{theorem}

\begin{corollary} \label{cor:enough-eigenvalues-implies-diagonalizable} \hlt{Enough Eigenvalues Imply Diagonalizability}\\
If $\mathcal{V}$ is finite-dimensional and $T\in\LL(\mathcal{V})$ has $\dim(\mathcal{V})$ distinct eigenvalues, then $T$ is diagonalizable.
\begin{proof}
Let $n=\dim(\mathcal{V})$, and let $\lambda_1,\ldots,\lambda_n$ be distinct eigenvalues of $T$. For each $j$, choose an eigenvector $\mathbf{v}_j$ corresponding to $\lambda_j$. By Theorem \ref{thm:distinct-eigenvalues-independent-eigenvectors}, $(\mathbf{v}_1,\ldots,\mathbf{v}_n)$ is linearly independent. Since this list has length $n$, Proposition \ref{prop:full-dimension-basis-criteria} makes it a basis of $\mathcal{V}$. Thus $\mathcal{V}$ has an eigenbasis, so $T$ is diagonalizable by Theorem \ref{thm:conditions-equivalent-to-diagonalizability}.\\
\hlt{Structural Commentary}: One eigenvector from each distinct eigenvalue already gives a full-size independent list.\\
\hlt{Key Insight}: Once the independent eigenvectors reach $\dim(\mathcal{V})$, they must form a basis.
\end{proof}
\end{corollary}

\begin{theorem} \label{thm:commutant-of-diagonalizable-matrix} \hlt{Commutant of a Diagonalizable Matrix}\\
Let $A\in\operatorname{Mat}(n,n,\mathbb{C})$ be diagonalizable. Suppose its distinct eigenvalues are $\mu_1,\ldots,\mu_d$, and set $n_i=\dim E_{\mu_i}(A)$.\\
Write $A=S(\mu_1I_{n_1}\oplus\cdots\oplus\mu_dI_{n_d})S^{-1}$. Then $\operatorname{Comm}(A)=\{S(X_1\oplus\cdots\oplus X_d)S^{-1}:X_i\in\operatorname{Mat}(n_i,n_i,\mathbb{C})\}$.\\
In particular, $\dim\operatorname{Comm}(A)=n_1^2+\cdots+n_d^2$.
\begin{proof}
Let $\Lambda=\mu_1I_{n_1}\oplus\cdots\oplus\mu_dI_{n_d}$. Matrix $B$ commutes with $A$ if and only if $S^{-1}BS$ commutes with $\Lambda$.\\
By Lemma \ref{lem:block-diagonal-commutant-distinct-scalars}, this happens exactly when $S^{-1}BS=X_1\oplus\cdots\oplus X_d$ with $X_i\in\operatorname{Mat}(n_i,n_i,\mathbb{C})$.\\
Dimension count follows because conjugation by $S$ is a vector-space isomorphism and $\dim\operatorname{Mat}(n_i,n_i,\mathbb{C})=n_i^2$.\\
\hlt{Structural Commentary}: Diagonalizing $A$ reduces commutation to respecting equal-eigenvalue blocks.\\
\hlt{Key Insight}: Distinct eigenvalue blocks cannot mix under a commuting matrix.
\end{proof}
\end{theorem}

\begin{theorem} \label{thm:simultaneous-schur-triangularization-commuting-family} \hlt{Simultaneous Schur Triangularisation}\\
Let $\emptyset\neq\mathcal{F}\subseteq\operatorname{Mat}(n,n,\mathbb{C})$ be a commuting family.\\
Then there is a unitary $U\in\operatorname{Mat}(n,n,\mathbb{C})$ such that $U^*AU$ is upper triangular for every $A\in\mathcal{F}$.\\
If every matrix in $\mathcal{F}$ is real and has only real eigenvalues, then $U$ can be chosen real orthogonal.
\begin{proof}
Induct on $n$. The case $n=1$ is immediate.\\
For $n>1$, Theorem \ref{thm:commuting-family-common-eigenvector} gives a common eigenvector; normalize it to a unit vector $\mathbf{x}$ and extend to an orthonormal basis. Let $V=[\mathbf{x}\ V_2]$ be the corresponding unitary matrix.\\
For every $A\in\mathcal{F}$, since $A\mathbf{x}=\lambda_A\mathbf{x}$, we have $V^*AV=
\begin{bmatrix}
\lambda_A&*\\
0&A'
\end{bmatrix}$.\\
If $A,B\in\mathcal{F}$, then $AB=BA$. The lower-right block of $V^*ABV=(V^*AV)(V^*BV)$ is $A'B'$, while the lower-right block of $V^*BAV$ is $B'A'$. Hence $A'B'=B'A'$, so the compressed family is still commuting.\\
By induction, choose a unitary $W$ that simultaneously upper triangularises all compressed blocks.\\
Then $U=V(1\oplus W)$ is unitary and $U^*AU$ is upper triangular for every $A\in\mathcal{F}$.\\
For the real case, run the same induction over $\mathbb{R}$. A common real eigenvector exists as follows. If all matrices are scalar, any real unit vector works. Otherwise choose a nonscalar $A_0\in\mathcal{F}$.\\
Since $A_0$ has only real eigenvalues, some real eigenspace $\mathcal{U}=E_{\lambda}(A_0)$ is nonzero; since $A_0$ is not scalar, $\mathcal{U}\neq\mathbb{R}^n$.\\
Commutation gives $(A_0-\lambda I)B\mathbf{u}=B(A_0-\lambda I)\mathbf{u}=\mathbf{0}$ for $B\in\mathcal{F}$ and $\mathbf{u}\in\mathcal{U}$, so $\mathcal{U}$ is invariant under every $B$.\\
The restricted family on $\mathcal{U}$ is commuting and still has only real eigenvalues: in a basis adapted to $\mathcal{U}$ each $B$ has block form $\begin{bmatrix}B|_{\mathcal{U}}&*\\0&*\end{bmatrix}$, so the characteristic polynomial of $B|_{\mathcal{U}}$ divides that of $B$. Induction on $\dim(\mathcal{U})$ gives a common real eigenvector. Extending by real orthonormal bases gives a real orthogonal $U$.\\
\hlt{Structural Commentary}: A common eigenline creates the first block; commutation descends to the compression.\\
\hlt{Key Insight}: Simultaneous triangularisation is Schur induction with a common eigenvector at every step.
\end{proof}
\end{theorem}

\begin{figure}[H]
  \centering
  \begin{tikzpicture}[
  >=Stealth,
  every node/.style={font=\small},
  box/.style={draw=gray!55, fill=gray!5, rounded corners=2pt, align=center, inner xsep=5pt, inner ysep=4pt},
  hi/.style={draw=blue!60!black, fill=blue!7, rounded corners=2pt, align=center, inner xsep=5pt, inner ysep=4pt},
  arr/.style={->, line width=0.75pt, blue!70!black}
]
  \node[hi] (common) at (0,0)
    {$\mathbb{C}\mathbf{x}$};
  \node[box] (block) at (3.18,0)
    {$V^*AV=\begin{bmatrix}\lambda_A&*\\0&A'\end{bmatrix}$};
  \node[box] (comm) at (6.60,0)
    {$A'B'=B'A'$};

  \draw[arr] (common.east) -- (block.west);
  \draw[arr] (block.east) -- (comm.west);
\end{tikzpicture}

  \caption{A shared eigenline gives a simultaneous first block; lower-right compressions still commute.}
  \label{fig:simultaneous-triangularization-commuting}
\end{figure}

\begin{corollary} \label{cor:commuting-sum-product-eigenvalues} \hlt{Eigenvalues of Sums and Products of Commuting Matrices}\\
Let $A,B\in\operatorname{Mat}(n,n,\mathbb{C})$ and $AB=BA$. Their eigenvalues can be ordered as $\lambda_1,\ldots,\lambda_n$ and $\mu_1,\ldots,\mu_n$ so that $A+B$ has eigenvalues $\lambda_i+\mu_i$ and $AB$ has eigenvalues $\lambda_i\mu_i$, counted with multiplicity.
\begin{proof}
Choose simultaneous Schur forms $T=U^*AU$ and $R=U^*BU$. Set $\lambda_i=t_{ii}$ and $\mu_i=r_{ii}$. Then $T+R$ and $TR$ are upper triangular with diagonals $\lambda_i+\mu_i$ and $\lambda_i\mu_i$, and are similar to $A+B$ and $AB$.\\
\hlt{Structural Commentary}: Simultaneous Schur form aligns both spectra on one diagonal.\\
\hlt{Key Insight}: In common triangular coordinates, diagonals add and multiply.
\end{proof}
\end{corollary}
